% Generated draft: requires mathematical and source review.
\documentclass[11pt]{article}
\usepackage[margin=25mm]{geometry}
\usepackage{fontspec}
\setmainfont{Latin Modern Roman}
\usepackage{amsmath,amssymb,mathtools,unicode-math}
\setmathfont{Latin Modern Math}
\usepackage{enumitem,needspace,xurl,hyperref}
\setlength{\parindent}{0pt}
\setlength{\parskip}{5pt}
\setlength{\emergencystretch}{3em}
\newcommand{\N}{\mathbb{N}}
\newcommand{\Z}{\mathbb{Z}}
\newcommand{\Q}{\mathbb{Q}}
\newcommand{\R}{\mathbb{R}}
\newcommand{\C}{\mathbb{C}}
\newcommand{\F}{\mathbb{F}}
\newcommand{\A}{\mathbb{A}}
\newcommand{\PP}{\mathbb{P}}
\newcommand{\E}{\mathbb{E}}
\newcommand{\eps}{\varepsilon}
\newcommand{\dd}{\,\mathrm{d}}
\begin{document}
\section*{UnsolvedMath ID: 700238}
\textbf{Source problem code:} catalogue entry
\subsection{Definitions and mathematical statement}
For a topological group \(H\), define its group of topological automorphisms by
\[
\operatorname{Aut}_{\mathrm{top}}(H)
=\{\alpha:H\to H\mid \alpha\text{ is a group isomorphism and both }\alpha\text{ and }\alpha^{-1}\text{ are continuous}\},
\]
with composition as the group operation. The word ``full'' means that all such automorphisms are included.

Let \(\mathcal T\) denote the class of topological groups permitted as realizing groups. For each of the target classes
\[
\mathcal C=\{\text{all groups}\},\qquad
\mathcal C=\{\text{finite groups}\},\qquad
\mathcal C=\{\text{abelian groups}\},
\]
characterize the groups \(G\in\mathcal C\) for which
\[
\exists H\in\mathcal T\quad G\cong\operatorname{Aut}_{\mathrm{top}}(H).
\]
Here \(\cong\) denotes abstract group isomorphism. Thus the requested description is a necessary and sufficient characterization of the groups arising as full automorphism groups of topological groups, with finite and abelian target groups singled out in the source.

The formulation remains qualified: the supplied record does not specify whether the realizing topological groups must be Hausdorff or satisfy another separation condition, so the intended class \(\mathcal T\) is unresolved. It also does not specify a topology on the automorphism group or on the target group. The display therefore gives the abstract-group interpretation; a classification up to topological group isomorphism would require additional data. The parenthetical phrase ``finite groups, abelian groups'' is retained as specifying target subclasses, although its precise intended relationship to the unrestricted classification is not explained in the record.
\subsection{Short English statement}
Describe which groups occur as the full groups of continuous automorphisms with continuous inverses of topological groups, including the finite and abelian target cases. The permitted topological groups and whether realization is meant abstractly or topologically require clarification.
\subsection{Sources}
\begin{itemize}
\item Local UnsolvedMath registry, permanent ID \(700238\), exact source problem code \texttt{catalogue entry}, title ``Kourovka Notebook Problem 4.30.'' The supplied registry labels the problem open; this status is not independently verified here.
\item M. I. Kargapolov, problem proposer, 1973. The supplied attribution identifies the original problem as Problem 4.30 in \emph{The Kourovka Notebook (Unsolved Problems in Group Theory)}, 4th issue (1973).
\item The supplied source identification also names the 21st issue edition (2026), \texttt{arXiv:1401.0300v48}. The collection description attributes editorial responsibility to E. I. Khukhro and V. D. Mazurov. No source text beyond the supplied record was independently consulted.
\end{itemize}
\end{document}
